\section{Largest Prime Factor --- Problem 3}
\label{prob:3}
\noindent\textbf{Source:} \url{https://github.com/KO0I/projectEuler/blob/reorganize/solutions-by-language/solutions/c/largest_prime_factor_3.c}

\noindent\textbf{Assessment:} Core method implemented.

\subsection*{Mathematical goal}
The goal is to determine the largest prime factor of \(600851475143\). Removing factors one by one isolates the greatest factor remaining at the end.

\subsection*{How the C code works}
The function \verb|factorize| starts with candidate divisor \verb|d=2|. If \verb|n%d==0|, it divides \verb|n| by \verb|d| and tries the same divisor again, removing repeated factors. Otherwise it moves to 3 and then successive odd candidates. Once the loop condition \verb|d<n| becomes false, it prints \verb|d|.

\subsection*{Correctness and limitations}
This factor-stripping approach yields the expected factor \texttt{...6857} for the supplied value. \verb|unsigned long| needs to be at least 64 bits to represent the input portably. Testing candidates only up to the square root of the remaining number could improve efficiency.


\subsection*{Correction to historical C --- Problem 3}
{\color{red}\textbf{Correction:} The corrected trial-division routine uses a 64-bit integer and stops factoring once a candidate divisor exceeds the square root of the remaining number.}
